\clearpage
\VerdierPhysicalPageAnchor{0114}%
\begingroup
\small
\setlength{\abovedisplayskip}{4pt plus 2pt minus 2pt}%
\setlength{\belowdisplayskip}{4pt plus 2pt minus 2pt}%

\VerdierSemanticAnchor{ch2-1-2}
\medskip
\noindent\textbf{1.2. 三角范畴的初步性质。}

\VerdierSemanticAnchor{ch2-1-2-1}
\medskip
\noindent\textbf{命题 1.2.1.}\quad
设 $\mathcal D$ 为一个三角范畴。对于 $\mathcal D$ 的每个对象 $X$，函子
$\operatorname{Hom}_{\mathcal D}(X,\mathord\cdot)$ 与
$\operatorname{Hom}_{\mathcal D}(\mathord\cdot,X)$ 都是取值于阿贝尔群范畴的
上同调函子。

\VerdierSemanticAnchor{ch2-1-2-1-hom-functors}
\noindent\emph{证明。}\quad
先证明函子 $\operatorname{Hom}_{\mathcal D}(X,\mathord\cdot)$ 是上同调函子
就够了。另一个断言随后可由取反范畴 $\mathcal D^\circ$ (1.1.7) 得到。
因此，设

\VerdierSemanticAnchor{ch2-1-2-1-source-triangle}
\[
  X'\xrightarrow{u'}Y'\xrightarrow{v'}Z'\xrightarrow{w'}X'[1]
\]
为一个特异三角。只需证明在图式

\VerdierSemanticAnchor{ch2-1-2-1-1}
\begin{equation}
  \resizebox{0.98\linewidth}{!}{$
    \operatorname{Hom}_{\mathcal D}(X,X')
    \xrightarrow{\operatorname{Hom}_{\mathcal D}(X,u')}
    \operatorname{Hom}_{\mathcal D}(X,Y')
    \xrightarrow{\operatorname{Hom}_{\mathcal D}(X,v')}
    \operatorname{Hom}_{\mathcal D}(X,Z')
    \xrightarrow{\operatorname{Hom}_{\mathcal D}(X,w')}
    \operatorname{Hom}_{\mathcal D}(X,X'[1])
  $}
  \tag{1.2.1.1}
\end{equation}

\VerdierSemanticAnchor{ch2-1-2-1-exactness-reduction}
中有
\[
  \operatorname{Hom}_{\mathcal D}(X,v')
  \operatorname{Hom}_{\mathcal D}(X,u')=0
\]
并且该图式在 $\operatorname{Hom}_{\mathcal D}(X,Y')$ 处正合。事实上，公理
TRII 随后会在 $\operatorname{Hom}_{\mathcal D}(X,Z')$ 处给出同一结论。
设 $f\colon X\longrightarrow X'$。由 TRI，三角

\VerdierSemanticAnchor{ch2-1-2-1-identity-triangle}
\[
  X\xrightarrow{\operatorname{id}_X}X\longrightarrow0\longrightarrow X[1]
\]
是特异三角；由 TRII 与 TRIII，存在态射 $g\colon X\longrightarrow Y'$，使下图
交换：

\VerdierSemanticAnchor{ch2-1-2-1-2}
\begin{equation}
  \begin{tikzcd}[column sep=3.2em,row sep=2.5em]
    X \arrow[r,"{\operatorname{id}_X}"] \arrow[d,"f"']
      & X \arrow[r] \arrow[d,"g"']
      & 0 \arrow[r] \arrow[d]
      & X[1] \arrow[d,"{f[1]}"] \\
    X' \arrow[r,"u'"']
      & Y' \arrow[r,"v'"']
      & Z' \arrow[r,"w'"']
      & X'[1]
  \end{tikzcd}
  \tag{1.2.1.2}
\end{equation}

\VerdierSemanticAnchor{ch2-1-2-1-zero-composite}
由此得到 $u'f=g$ 以及 $v'u'f=v'g=0$，从而证明
$\operatorname{Hom}_{\mathcal D}(X,v')
\operatorname{Hom}_{\mathcal D}(X,u')=0$。

\VerdierSemanticAnchor{ch2-1-2-1-lifting}
现设 $g\colon X\longrightarrow Y'$ 为满足 $v'g=0$ 的态射。由 TRII 与
TRIII，存在态射 $f\colon X\longrightarrow X'$，使图 (1.2.1.2) 交换。
于是 $g=u'f$，因此列 (1.2.1.1) 在
$\operatorname{Hom}_{\mathcal D}(X,Y')$ 处正合。

\VerdierSemanticAnchor{ch2-1-2-2}
\medskip
\noindent\textbf{推论 1.2.2.}\quad
在一个特异三角
$X'\xrightarrow{u'}Y'\xrightarrow{v'}Z'\xrightarrow{w'}X'[1]$ 中，
任意两个相继态射的复合为零。

\VerdierSemanticAnchor{ch2-1-2-2-proof}
\noindent\emph{证明。}\quad
由 TRII，只需证明 $v'u'=0$。为此，在正合列 (1.2.1.1) 中取
$X=X'$，并考察
$\operatorname{id}_{X'}\in\operatorname{Hom}_{\mathcal D}(X',X')$ 的
相继像即可。

\VerdierSemanticAnchor{ch2-1-2-3}
\medskip
\noindent\textbf{推论 1.2.3.}\quad
设 $(X,Y,Z,u,v,w)$ 与 $(X',Y',Z',u',v',w')$ 为 $\mathcal D$ 中两个
特异三角，$(f_1,f_2,f_3)$ 为三角的一个态射。如果三个 $f_i$ 中的两个
（$1\leq i\leq3$）是同构，则第三个也是同构。
\endgroup
